\section{Introduction}
\label{sec:introduction}

\subsection{Motivation}
\label{sec:introduction-motivation}

Modern main-memory transaction processing systems exploit many-core machines by
removing centralized bottlenecks from concurrency control.  Timestamp-based
engines, including ERMIA-style protocols, already maintain a commit timestamp
that determines the serialization order of committed transactions.  As a result,
the critical path of transaction execution is no longer determined only by
validation or conflict management.  When crash durability is required, the
write-ahead logging (WAL) protocol can become the dominant scalability limit.

A conventional centralized WAL provides a simple and safe durability order, but
it serializes log insertion and durable acknowledgment through a shared log
stream.  Parallel WAL (P-WAL) removes this obvious serialization point by
partitioning log records across multiple WAL shards.  However, once the shared
WAL mutex is removed, the system exposes the next durability bottlenecks:
storage synchronization, durable-point notification, and waiting for a global
durable prefix.  These costs are particularly important for IPSJ Transactions on
Databases (TOD)-style database systems research, where the durability protocol
must be evaluated as part of the end-to-end transaction processing design rather
than as an isolated I/O optimization.

\subsection{Problem}
\label{sec:introduction-problem}

The key problem is that existing P-WAL designs often reintroduce a conservative
global ordering mechanism at the durability layer.  A transaction may have an
ERMIA commit timestamp that already defines its serialization position, but the
WAL layer separately allocates a global log sequence number (LSN) and acknowledges
durability only when a global durable prefix reaches that LSN.  This design is
safe, but it couples the acknowledgment of unrelated transactions.  A slow WAL
shard can delay transactions on other shards even when they do not depend on any
log record written by the slow shard.

The opposite design point is local-only acknowledgment: acknowledge a transaction
as soon as its own WAL shard is durable.  This removes the global-prefix wait and
can improve throughput and latency, but it is unsafe in general.  If transaction
\(T\) reads data produced by transaction \(U\), then acknowledging \(T\) before
\(U\)'s log record is durable can make recovery replay \(T\) without the state on
which \(T\) depended.  Therefore, the durability protocol must avoid both
unnecessary global-prefix waiting and unsafe local-only acknowledgment.

\subsection{Approach}
\label{sec:introduction-approach}

This paper proposes CStamp-PWAL, a dependency-closed parallel WAL protocol for
timestamp-based transaction processing.  The first idea is to reuse the commit
timestamp as a logical LSN.  The WAL layer still maintains physical positions
within each shard, but it does not need to allocate a separate global LSN for the
recovery order.  This removes a duplicated ordering mechanism between
concurrency control and durability.

The second idea is dependency-closed durable acknowledgment.  Instead of waiting
for a global durable prefix, CStamp-PWAL acknowledges a transaction only after
its own log record and the durable frontier of the transactions on which it
depends have become durable.  This condition preserves the dependency closure
required for safe recovery while allowing independent WAL shards to progress
without waiting for unrelated slow shards.

The third idea is to separate transaction execution from durable acknowledgment.
Workers execute transactions, validate them, assign commit timestamps, and enqueue
log records.  Flushers batch and persist shard-local log records.  Committers
observe durable-frontier advances and acknowledge transactions whose dependency
conditions are satisfied.  This pipeline keeps worker threads from blocking
directly on \texttt{fdatasync} or on conservative durable-prefix waits.

\subsection{Organization}
\label{sec:introduction-organization}

The rest of this paper is organized as follows.  Section~2 reviews WAL,
parallel WAL, and dependency requirements for crash recovery.  Section~3
describes the design of CStamp-PWAL, including commit timestamps as logical LSNs,
dependency-frontier construction, and the worker/flusher/committer pipeline.
Section~4 presents the experimental methodology and evaluates the throughput,
backlog, and tail-latency effects of CStamp-PWAL.  Section~5 discusses related
work, and Section~6 concludes the paper.
